\subsection{Insertion sort}

\textit{Insertion sort} resuelve el problema de ordenamiento, definido en la sección \ref{subsec:problema}.

Si el problema es ordenar todo el arreglo, una versión más pequeña es ordenar un pedazo del arreglo. Siguiendo el paradigma, imaginemos que ese problema más pequeño ya está resuelto: que todos los elementos del arreglo menos el último ya están ordenados. El paso que falta es obvio: tomar ese último elemento e insertarlo en la posición correcta, comparando hacia atrás.

Esa es la idea que da origen a \concept{insertion sort}: para cada elemento del arreglo, se revisa si está ubicado correctamente respecto a los anteriores: si no lo está, se mueve hacia atrás hasta que lo esté; una vez lo está, se continúa con el siguiente elemento.

\begin{pseudocode}
def insertion_sort(nums: arreglo):
    |\bxl{s}|i = 1|\bxr{s}|

    while i < len(nums):
        elem = nums[i] 
        j = i-1
        |\br{anti}|while j >= 0 and elem < nums[j]:
            nums[j+1] = nums[j]
            j -= 1|\br{antf}|
        
        |\bxl{in}|nums[j+1] = elem|\bxr{in}|
        i += 1
\end{pseudocode}
\begin{annotations}
  \codebox[xshift=4.5cm, width=10cm]{s}{Se inicia revisando el segundo elemento, para poder compararlo con algún anterior.}
  \codebrace[xshift=4cm, width=9cm]{anti}{antf}{Si \inline{elem} es menor que el anterior y debe moverse hacia atrás, se mueve el anterior una posición adelante para hacerle espacio. Se siguen corriendo los anteriores hasta que ya no deba moverse \inline{elem} hacia atrás (o ya no hayan anteriores).}
  \codebox[width=10cm]{in}{Se inserta \inline{elem} en la posición que quedó libre tras mover los elementos más grandes que él hacia adelante. Es \inline{j+1} y no \inline{j} porque la última iteración del ciclo termina con \inline{j -= 1}.}
\end{annotations}

En cada iteración se acomoda un elemento, reduciendo el problema en una constante: uno.

\subsubsection{Complejidad}

El algoritmo presenta un ciclo anidado dentro de otro. El ciclo externo siempre recorre el arreglo completo, pero la cantidad de veces que ejecuta el cuerpo del ciclo interno depende de qué tanto haya que retroceder cada elemento, es decir, de qué tan ordenado esté el arreglo original.

Si el arreglo está ordenado, no hay que retroceder ningún elemento y el cuerpo del ciclo interno nunca ejecuta. Ese es el mejor caso: el costo computacional está dado solo por el ciclo externo, que revisa todos los elementos. Sea \(T_\text{mejor}\) la función para ese costo computacional, \(T_\text{mejor} \in \Theta(n)\).

El peor caso es cuando el ciclo interno tiene que retroceder cada elemento hasta el principio, es decir, cuando el arreglo está en orden inverso. En ese caso, para cada elemento \(i\), el ciclo interno debe retrocederlo \(i-1\) veces; la cantidad de operaciones es entonces:
\[\sum_{i=2}^n (i-1) = 1+2+\cdots+(n-1) = \frac{n^2-n}{2}.\]
Eso, aunado a las operaciones del ciclo externo, da la siguiente función para el costo computacional, donde \(c\) denota el tiempo que demoran las operaciones constantes:
\[T_\text{peor}(n) = cn + c\frac{n^2-n}{2}\]
Ergo, \(T_\text{peor} \in \Theta(n^2)\). Convencionalmente, se dice que insertion sort es \(\Theta(n^2)\), entendiendo que se hace referencia específicamente al peor caso.

\begin{tip}
Para resolver la sumatoria se empleó la fórmula de Gauss, cuyo argumento subyacente es bastante simple y se expone en el apéndice \ref{apx:sumatorias}.
\end{tip}